% ---
% id: meat-core/frank-wolfe/cor:fo-reference
% kind: cor
% title: "The first-order bound on the reference companion"
% summary: >-
%   At any points, such as those of a search on the snapped declaration or
%   of one continued on the companion, the bracket formed with the reference
%   companion's objective, gradient, entries and linear minima bounds the
%   companion's value from below; its key length is the reference model's.
% prominence: main
% status: accepted
% uses: [thm:firstorder, cor:fo-keylength, def:relaxed-feasible-set, def:first-order-setting, meat-core/prop:reference-companion, meat-core/def:admissible-declaration, meat-core/def:admissible-acceptance-set, meat-core/lem:branch-support, meat-core/def:joint-sectors]
% mentions: [req:bracket-on-reference, rem:fo-arithmetic, meat-core/rem:declared-measurement]
% assumes: []
% refs: []
% provenance: newly-derived
% lean:
% notes: "Read for the companion of meat-core/prop:reference-companion, which may replace several measured factors, so that (ii) and its proof speak of the replaced links."
% ---

Theorem~\ref{thm:firstorder} asks nothing about where its points came from:
they may be any finite set of points of the relaxed feasible set at which its
hypotheses hold. A run may therefore search on the declaration $\Decl$
compiled from the snapped measurement and form the bracket \eqref{eq:hlb} on
the compiled reference companion $\Decl^{\mathrm{ref}}$ at the points the
search produced. The snap then only steers the search, and its displacement
plays no part in the bound
(Remark~\ref{rem:declared-measurement}). The corollary is the theorem read
for $\Decl^{\mathrm{ref}}$, with the hypotheses that carry over from $\Decl$
named; what a run evaluates, and how, is
Requirement~\ref{req:bracket-on-reference}, and the arithmetic of that
evaluation is Remark~\ref{rem:fo-arithmetic}.

\begin{corollary}\label{cor:fo-reference}
Let $\Decl$ and $\Decl^{\mathrm{ref}}$ be as in
Proposition~\ref{prop:reference-companion}, let $\alpha\in(1,\infty)$, and let
$\acSet$ be an \tref{admissible acceptance set}
(Definition~\ref{def:admissible-acceptance-set}), the one of the run, read
for both. Let $\varepsilon\ge0$ and let $x_t=(\mathbf q_t,J_t)$,
$t\in\Itset$, be a finite non-empty set of points, produced by any
computation, a search on $\Decl$ among them, such that every $x_t$ lies in
$\acSet\times\Feas^{J}(\varepsilon)$ of $\Decl^{\mathrm{ref}}$
(Definition~\ref{def:relaxed-feasible-set}) and satisfies (H-POS), (H-SUPP)
and (H-NOPERT) of Theorem~\ref{thm:firstorder} for the objective, the \tref[branch map]{branch
maps} and the gradient of $\Decl^{\mathrm{ref}}$, and let $(v_t,S_t)$
satisfy \eqref{eq:lmo} at the gradient of $\Decl^{\mathrm{ref}}$ at $x_t$,
its feasible-side minimum over the exact $\Feas^{J}$ of $\Decl^{\mathrm{ref}}$
(Definition~\ref{def:first-order-setting}), or let, at some or all $t$,
lower bounds on one or both of those minima stand in for the pairings as
Theorem~\ref{thm:firstorder}\,(iv) lets them, \eqref{eq:lmo-bounds} read
at that gradient and over those sets. Write
$\hlb[\Decl^{\mathrm{ref}}]$ for \eqref{eq:hlb} so formed. Then:
\begin{enumerate}[label=(\roman*),leftmargin=2.4em]
\item in exact real arithmetic,
  $\max\{0,\hlb[\Decl^{\mathrm{ref}}]\}\le\hlow[\Decl^{\mathrm{ref}}]$, and
  $\hlow[\Decl^{\mathrm{ref}}]\le\hlow[\Decl^{\mathrm{ref}}_{\mathrm{ex}}]$
  for the companion with its unrelaxed entries;
\item the Choi variables of $\Decl$ and $\Decl^{\mathrm{ref}}$ are lists over
  the same \tref[joint sector]{joint sectors} with the same block dimensions,
  so every point of the one program is a point of the other; where no link
  the companion replaces carries entries, as on the encoded route,
  $\acSet\times\Feas^{J}(\varepsilon)$ is the same set for both; and a point with $J_{t,s}\succ0$ for every sector
  $s$ satisfies (H-SUPP) for $\Decl^{\mathrm{ref}}$, whatever its measured
  family;
\item where the hypotheses of Corollary~\ref{cor:fo-keylength} hold for
  $\Decl^{\mathrm{ref}}$, its conclusions hold with
  $\hlb[\Decl^{\mathrm{ref}}]$ in \eqref{eq:keylength-fo}: the key length so
  formed is secret and correct for the protocol whose declaration is
  $\Decl^{\mathrm{ref}}$.
\end{enumerate}
\end{corollary}

\begin{proof}
(i) By Proposition~\ref{prop:reference-companion}\,(ii),
$\Decl^{\mathrm{ref}}$ is a tuple of the form of
Definition~\ref{def:admissible-declaration} that satisfies its conditions
(i)--(vi) and (viii). Where it satisfies (vii) as well it is an admissible
declaration, and Theorem~\ref{thm:firstorder}\,(iiib) applies to it with
the given points and pairs, with its part~(iv) where lower bounds stand in,
whose hypotheses are exactly those listed: membership in the relaxed set of
$\Decl^{\mathrm{ref}}$, (H-POS), (H-SUPP) and (H-NOPERT) for its objects,
and \eqref{eq:lmo}, or \eqref{eq:lmo-bounds} where a lower bound stands in,
over the exact sets of $\Decl^{\mathrm{ref}}$. The theorem reads of the
points nothing but these hypotheses, so it holds whatever computation
produced them. Where (vii) fails, Theorem~\ref{thm:firstorder}\,(v)
gives the same conclusion, $\hlow[\Decl^{\mathrm{ref}}]=+\infty$. The second
inequality is Proposition~\ref{prop:reference-companion}\,(iii).

(ii) The Choi variable of a declaration in the general format is a list of
positive semidefinite blocks over its joint sectors, of dimension the
product of the Alice sector's and the measured sector's
(Definition~\ref{def:joint-sectors}); these are shared by
Proposition~\ref{prop:reference-companion}\,(i). The relaxed set of
Definition~\ref{def:relaxed-feasible-set} is cut out by per-block positivity,
the coupling condition and the declared entries; the first two are shared,
and the entries differ only in those of the replaced links
(Proposition~\ref{prop:reference-companion}\,(i)), so where none of those
links has any the two sets coincide. By
Lemma~\ref{lem:branch-support}\,(i), $J_{t,s}\succ0$ gives every
key-active branch operator its maximal support, for any declaration, which
is (H-SUPP).

(iii) Corollary~\ref{cor:fo-keylength} holds for any $\hlb$ of
Theorem~\ref{thm:firstorder}\,(ii) or (iiib), with the lower bounds of its
part~(iv) where they stand in, under its hypotheses, and (i) gives one for
$\Decl^{\mathrm{ref}}$.
\end{proof}
